\chapter{Esophageal pH Test}
\section*{What Is This Test?} Esophageal pH monitoring measures acid exposure in the esophagus over an extended period to evaluate reflux.
\section*{Alternative Names} Ambulatory pH monitoring; 24-hour pH test; pH-impedance monitoring.
\section*{Principle} A thin catheter or wireless capsule records acid episodes; impedance can also detect non-acid reflux.
\section*{Why This Test Is Ordered} It investigates persistent heartburn, regurgitation, cough, chest symptoms, or suspected reflux when diagnosis is uncertain or treatment is ineffective.
\section*{Primary vs Secondary Test} It is a secondary physiologic test, often selected after history and endoscopy; manometry may be needed before catheter placement.
\section*{How This Test Is Performed} A catheter is passed through the nose or a capsule is attached during endoscopy. The patient records meals, position, symptoms, and medicines.
\section*{What Preparation Is Needed} Acid-suppressing medicines may need adjustment only under clinician direction. Follow fasting and activity instructions.
\section*{Understanding Results} The report measures acid exposure and correlation with symptoms. A normal result makes pathologic acid reflux less likely but does not explain every symptom.
\section*{Follow-up Tests} Endoscopy, manometry, swallow imaging, gastric evaluation, or specialist consultation may follow.
\section*{References} \begin{enumerate}\item MedlinePlus. Esophageal pH Test.\item American College of Gastroenterology. GERD diagnostic guidance.\end{enumerate}
\clearpage\section*{Understanding Results and Follow-up}
The laboratory report should identify the specimen, method, units, reference interval or decision limit, and whether the result is positive, negative, abnormal, or inconclusive. Reference intervals vary with age, sex, pregnancy, specimen type, assay, and laboratory; a result outside a range is not automatically a diagnosis, and a result within a range does not always exclude disease. Discuss unexpected results with the ordering clinician, who can decide whether repeat or confirmatory testing, treatment, or referral is appropriate. Do not change medicines or delay urgent care based on an isolated result.
\section*{Regulatory Status and Authorization Date}This chapter describes a test category, not a specific commercial product. FDA, CDSCO, EU IVDR, and MHRA approval or registration dates therefore cannot be assigned without the assay manufacturer, product name, instrument, and intended use. For laboratory-developed tests, an individual agency approval date may not exist. See the Preface for the interpretation of regulatory status.
\clearpage
